\subsection{Thảo luận và đánh giá}
\label{sec:discussion}

\subsubsection{Cơ chế đạt được gì so với công trình tham chiếu}
Cơ chế đề xuất tái hiện được \emph{định hướng và tính khả thi} của lõi phát hiện
tham chiếu DeSFAM: bộ phát hiện VAE\,+\,iForest xếp hạng bất thường tốt (AUC kiểm
định VAE $0{,}944$; AP tập hợp $0{,}990$) và giữ khả năng phân tách khi phục vụ
trực tiếp qua Tetragon. Điều \emph{chưa} đạt là các con số tuyệt đối ở chế độ phân
loại nhị phân---đặc biệt F1 $0{,}92$ và recall cao của công trình tham chiếu. Nguyên nhân khả dĩ:
(i) công trình tham chiếu
chọn ngưỡng tối ưu F1 hoặc đánh giá trên một phân chia tấn công/khối lượng khác;
(ii) khác biệt về phiên bản DongTing, cách cắt cửa sổ và tập syscall đầu vào;
(iii) công trình tham chiếu bổ sung Giai đoạn~3 (thực thi trong nhân) để đạt ``chặn 14/14 CVE''
---một chỉ số \emph{thực thi}, không phải chỉ số phân loại của riêng SyscallAD.

\subsubsection{Ý nghĩa của khoảng cách ngưỡng}
Việc AUC/AP cao nhưng F1 thấp tại ngưỡng cố định cho thấy giá trị của \emph{hiệu
chỉnh ngưỡng theo môi trường}. Ngưỡng bảo thủ chọn trên tập kiểm định ($0{,}81$
cho điểm tập hợp đã chuẩn hóa) tối đa hóa precision nên bỏ sót cửa sổ tấn công
``nhẹ''. Ngược lại, khi phục vụ trực tiếp, ngưỡng vận hành $T_0=0{,}25$ được hiệu
chỉnh theo phân phối quan sát giúp tách sạch hai lớp. Đây là một phát hiện thực
dụng: với phát hiện bất thường, hiệu chỉnh ngưỡng tại chỗ quan trọng ngang với
chất lượng mô hình.

\subsubsection{Lựa chọn thiết kế khác với công trình tham chiếu}
Chúng tôi dùng Tetragon~\cite{tetragon} làm cảm biến eBPF thay vì tự viết chương
trình eBPF/LSM. Điều này đẩy nhanh hiện thực hóa và bám hệ sinh thái Kubernetes
thực tế, nhưng cũng \emph{thu hẹp} bảng chữ cái syscall về 23 lời gọi an ninh---một
ràng buộc kéo theo phải lọc dữ liệu huấn luyện tương ứng
(Mục~\ref{sec:evaluation}). Việc thực thi/chặn trong nhân (Giai đoạn~3 đầy đủ của
DeSFAM) chưa được hiện
thực hóa; chuyên đề dừng ở phát hiện và cảnh báo.

\subsubsection{Hạn chế và đe dọa tính hợp lệ}
\begin{itemize}[leftmargin=1.6em]
  \item \textbf{Hợp lệ nội tại:} kịch bản tấn công trực tiếp là \emph{mô phỏng}
        chuỗi syscall đặc quyền, không phải khai thác CVE thực; do đó độ phân tách
        ở Bảng~\ref{tab:live} là dấu hiệu khả quan chứ chưa phải bằng chứng chặn
        khai thác thật.
  \item \textbf{Hợp lệ ngoại tại:} một khối lượng lành tính (MongoDB) và một không
        gian tên thử nghiệm chưa đại diện cho đa dạng khối lượng sản xuất; cần mở
        rộng (Nginx, Redis, Node.js\dots) như công trình tham chiếu.
  \item \textbf{Hợp lệ kết luận:} mất cân bằng lớp lớn (tỉ lệ tấn công cao trong
        tập kiểm tra cửa sổ) khiến AP cao dễ gây hiểu nhầm; chúng tôi báo cáo kèm
        AUC, precision, recall để tránh kết luận một chiều.
  \item \textbf{Hợp lệ cấu trúc:} đặc trưng thời gian bị tắt do DongTing không có
        dấu thời gian theo syscall; mô hình có thể thiếu một tín hiệu mà công trình
        tham chiếu tận dụng.
\end{itemize}

\subsubsection{Hướng cải thiện}
(1) Chọn ngưỡng tối ưu F1 và báo cáo đường cong precision--recall đầy đủ;
(2) hiện thực hóa Giai đoạn~3 (chặn bằng eBPF/LSM hoặc Tetragon Enforcement) để đo
``tỉ lệ chặn'' so sánh trực tiếp với công trình tham chiếu; (3) mở rộng tập khối lượng và dùng
khai thác CVE thực trong môi trường cách ly; (4) bổ sung đặc trưng thời gian khi
dùng cảm biến có dấu thời gian (Tetragon cung cấp), điều DongTing thiếu.
